\chapter{アプローチの検証：切断中の安定把持の作業}
% ============================================================
\label{ch:holding_application}
% ============================================================

% ------------------------------------------------------------
\section{問題設定と本章の位置づけ}
\label{sec:holding_problem_and_role}
% ------------------------------------------------------------

% ■ SICE2025 非対応（D論独自）
\subsection{本章の位置づけ}
\label{subsec:holding_position}

% --- 段落1 ---
本章は，本論文が提案する「センシング→モデリング→幾何・物理的分析→動作制御」という統一的な操作設計アプローチを，包丁による食材切断という強い外力と道具拘束が同時に作用するタスクへ適用した事例である．

% --- 段落2 ---
第3章で述べた3Dモデル取得基盤の上に，本章では食材表面上での把持点探索と包丁外力下での力学解析を積み上げ，双腕ロボットによる把持・切断動作の実行へと接続する．

% --- 段落3 ---
本タスクは，包丁から食材に加わる力とトルクに釣り合える食材の意図しない移動を防ぐ把持点集合を，食材の三次元形状と包丁軌道を入力として解析的に導出する問題として定式化される．


% ■ SICE2025 III-A (Problem formulation) + III-B (Method Flow) の概要
\subsection{提案手法の概要とアルゴリズムの流れ}
\label{subsec:holding_method_overview}

% --- 段落1 ---
提案手法は，食材表面の離散点群モデル$\Omega$，包丁軌道$\mathbb{T}$，および指の本数$n$を入力とし，切断中の全時刻を通じて最も安定な把持を実現する把持点集合$P_{\mathrm{fin}}$を出力するオフラインの探索アルゴリズムである．

% --- 段落2 ---
アルゴリズムは包丁軌道の各時刻ステップにおいて，(1) 現在の包丁姿勢に基づく把持候補の生成，(2) 幾何フィルタによる候補空間の削減，(3) 位置評価スコアの計算，(4) 力とトルクの推定，(5) 把持力候補の生成と動的評価スコアの計算，を順次実行する．

% --- 段落3 ---
全時刻の評価スコアを累積した後，最高スコアを得た把持点集合を最終出力とすることで，単一の包丁姿勢ではなく包丁軌道全体を考慮した把持計画を実現する点に本手法の特徴がある．


% ■ SICE2025 VI-A-2 (Surface Acquisition)
\subsection{点群情報の利用}
\label{subsec:holding_pointcloud}

% --- 段落1 ---
本手法では，まずRGB-Dカメラをロボットアーム先端に搭載し，食材の周囲複数視点から点群を取得する．

% --- 段落2 ---
取得した多視点点群は世界座標系へ変換・統合した後，移動最小二乗法（MLS）により平滑化し，さらに約150〜200点へダウンサンプリングする．

% --- 段落3 ---
このダウンサンプリング点数は，計算コストと評価スコアの関係を分析した上で経験的に決定した値であり，点数を増やしてもスコアはほとんど改善せず計算時間のみが増大することが確認されている．

% --- 段落4 ---
また，力とトルクの推定に必要な切断接触面$\Omega_{\mathrm{c1}},\Omega_{\mathrm{c2}}$は，取得点群から直接得ることができないため，包丁切断平面への点群投影とDelaunay三角分割により計算的に生成する．


% ------------------------------------------------------------
\section{分析：包丁外力下における把持安定性の力学的評価}
\label{sec:holding_analysis}
% ------------------------------------------------------------

% ■ SICE2025 III-B (Method Flow) + Algorithm 1
\subsection{アルゴリズムの全体流れ}
\label{subsec:holding_algorithm_flow}

% --- 段落1 ---
提案アルゴリズムの全体像をAlgorithm 1（Algorithm 1の参照ラベル）に示す．

% --- 段落2 ---
まず，食材表面$\Omega$の$K$個の離散点から$n$点を選ぶ全組合せとして全把持点集合$\mathbb{P}_{\mathrm{all}}$を生成し，累積スコアを格納する辞書$\mathbb{S}$を初期化する．全把持点集合の要素数は組合せ数$\binom{K}{n}$となる．

% --- 段落3 ---
メインループでは，包丁軌道$\mathbb{T}$の各時刻$(T,\bm{v})$に対し，現在の包丁姿勢$T$に基づいて食材表面を把持部$\Omega_{\mathrm{g}}$と切断接触面$\Omega_{\mathrm{c1}},\Omega_{\mathrm{c2}}$に分割する．その後，幾何フィルタ（\ref{subsec:holding_geometric_filter}節），位置評価スコア（\ref{subsec:holding_positional_score}節），力とトルクの推定（\ref{subsec:holding_knife_wrench}節），動的評価スコア（\ref{subsec:holding_dynamic_score}節）を順次計算し，各把持点集合のスコアを$\mathbb{S}$に加算する．各時刻で把持部$\Omega_{\mathrm{g}}$に含まれない把持点集合は無効とし，スコア$-\infty$を割り当てる．

% --- 段落4 ---
全時刻の走査が完了した後，$\mathbb{S}$において最大値をとる把持点集合$P_{\mathrm{fin}}$を最終出力とする．


% ■ SICE2025 IV (Search Space Reduction) + Algorithm 2
\subsection{幾何フィルタによる探索空間の削減}
\label{subsec:holding_geometric_filter}

% --- 段落1 ---
全把持点集合$\mathbb{P}_{\mathrm{all}}$の要素数は点群点数$K$に対して膨大となるため，後続の力学的評価に先立ち，低計算コストな幾何的制約のみを用いたフィルタリングにより探索空間を削減する（Algorithm 2）．

% --- 段落2 ---
幾何スコア$S_{\mathrm{geo}}$は，「指同士の距離」「指と包丁の距離」「指の高さ」の3要素の重み付き和として定義される．

% --- 段落3 ---
「指同士の距離」評価$E_{\mathrm{fin}}$は，二指が過度に接近して実質的に一指として機能することを防ぐために導入される．指間距離が閾値$b$を超えれば満点1となり，それ以下の場合は指数関数的に減衰する設計とした．

% --- 段落4 ---
「指と包丁の距離」評価$E_{\mathrm{knf}}$は，包丁の中心平面から各指先位置までの最短距離として定義され，切断動作中の手指と包丁の衝突を防止する役割を担う．

% --- 段落5 ---
「指の高さ」評価$E_{\mathrm{tbl}}$は，各指先の$z$座標の最小値として定義され，把持点がまな板に近すぎる候補を排除する．

% --- 段落6 ---
これら3つの評価値を$[0,1]$に正規化した後，重み$w_{\mathrm{fin}}, w_{\mathrm{knf}}, w_{\mathrm{tbl}}$により結合し，上位$\lfloor rN\rfloor$件の候補のみを後段へ出力する．本実験では出力比率$r=0.6$とした．なお，本フィルタでは指先が食材に接触した状態における拘束のみを考慮し，非接触状態から接触に至る指先の移動経路における自己衝突や環境衝突については扱わない．


% ■ SICE2025 V-A (Calculation of positional score) + Algorithm 3
\subsection{位置評価スコア}
\label{subsec:holding_positional_score}

% --- 段落1 ---
位置評価スコア$S_{\mathrm{pos}}$は，把持点集合の幾何的な配置が包丁から加わる力とトルクへの抵抗に適しているかを評価する指標であり，「方向」と「距離」の二要素から構成される（Algorithm 3）．

% --- 段落2 ---
「方向」評価$E_{\mathrm{pdir}}$は，把持点群の主成分方向ベクトル$\mathrm{PCA}(P)$と包丁平面法線$\hat{\bm{n}}_{\mathrm{k}}$の内積に基づいて定義される．主成分方向が包丁平面に平行であるほど高得点となり，これは多方向からの外力に対して有効に抵抗できる把持配置を優先する意図に基づく．

% --- 段落3 ---
「距離」評価$E_{\mathrm{pdis}}$は，包丁中心平面から各把持点までの距離の最小値に負号を付して定義される．包丁平面に近い位置で把持すればより小さな指先力でモーメントを打ち消せるという力学的直観を反映した設計である．

% --- 段落4 ---
両評価値を$[0,1]$に正規化し，重み$w_{\mathrm{pdir}}, w_{\mathrm{pdis}}$により結合することで$S_{\mathrm{pos}}$を得る．


% ■ SICE2025 V-B / III-B (Estimation of knife wrench)
\subsection{包丁が食材に与える力とトルクの推定}
\label{subsec:holding_knife_wrench}

% --- 段落1 ---
包丁が食材に及ぼす力学的負荷$\bm{t}_{\mathrm{k}}\in\mathbb{R}^6$は，Muら\cite{Mu2019RoboticCutting}の手法に基づき，食材の破壊靭性に起因する切断力と，包丁側面と食材接触面間に生じる摩擦力の二成分に分解して推定される．

% --- 段落2 ---
切断力$\bm{f}_{\mathrm{c}}$は，包丁刃線上の各位置$u$における速度方向$\hat{\bm{v}}(u)$と刃面法線$\hat{\bm{n}}_s(u)$の内積に食材の破壊靭性$\kappa$を乗じた微小切断力を，刃線全体にわたって積分することで求められる．また，食材の重心位置$\bm{g}$を用いて切断力が食材に与えるトルク$\bm{\tau}_{\mathrm{c}}$も併せて計算される．

% --- 段落3 ---
摩擦力$\bm{f}_{\mathrm{fr}}$は，切断接触面$\Omega_{\mathrm{c}j}$上の圧力分布が一様であるとの仮定のもと，クーロンの摩擦則に基づいて推定される．包丁速度$\bm{v}$を各接触面へ射影した速度成分$\bm{v}_{\mathrm{c}j}$を求め，摩擦係数$\mu$と接触面積を用いて積分する．

% --- 段落4 ---
以上の切断力と摩擦力を6次元ベクトルとして合成した$\bm{t}_{\mathrm{k}}=[(\bm{f}_{\mathrm{fr}}+\bm{f}_{\mathrm{c}})^{\mathrm{T}},\ (\bm{\tau}_{\mathrm{fr}}+\bm{\tau}_{\mathrm{c}})^{\mathrm{T}}]^{\mathrm{T}}$を，包丁が食材に与える力とトルクとして定義する．


% ■ SICE2025 V-B3 (Holding force & Finger force)
\subsection{把持力候補の効率的生成}
\label{subsec:holding_force_generation}

% --- 段落1 ---
把持点集合$P$が定まったとしても，包丁から食材に加わる力とトルク$\bm{t}_{\mathrm{k}}$と釣り合うべき把持力$\bm{f}\in\mathbb{R}^{3n}$には，把持行列$G$の零空間に起因する無限の冗長性が存在する．すなわち，$\bm{f}=G^{+}\bm{t}_{\mathrm{h}}+[I_{3n}-G^{+}G]\bm{k}$（$\bm{k}$は任意の$3n$次元ベクトル）と表され，$\bm{k}$の選び方により無数の把持力が得られる．ただし，$\bm{t}_{\mathrm{h}}$は$\bm{t}_{\mathrm{k}}$のうち把持力により打ち消すべき成分を抽出したものである．

% --- 段落2 ---
この冗長性に対処するため，本手法ではまず多数の把持力候補を生成し，その中から最適なものを選択する戦略をとる．

% --- 段落3 ---
最も単純な生成法は，任意ベクトル$\bm{k}$に乱数を割り当てる方法である．しかし，このアプローチでは生成された指先力の多くが摩擦円錐の外側を向く，すなわち物理的に妥当でない解となり，生成効率が極めて低い．

% --- 段落4 ---
そこで提案手法では，二段階の生成法を採用する．第一段階では，各指先力の方向を摩擦円錐内に限定した初期力$\bm{f}_{\mathrm{init}}$をランダムに生成する．第二段階では，この$\bm{f}_{\mathrm{init}}$と包丁から食材に加わる力とトルクとの不釣合い分を把持行列の擬似逆行列$G^{+}$により補正し，釣り合いを満たす把持力$\bm{f}$を得る．

% --- 段落5 ---
この提案手法により，摩擦円錐内に収まる妥当な把持力候補を，ランダム$\bm{k}$法と比較して平均約7.7倍の効率で生成できることが実験的に確認されている（\ref{subsec:holding_exp_force_generation}節参照）．それでも妥当な指先力が一つも生成できない把持点集合については，無効候補としてスコア$-\infty$を割り当てる．


% ■ SICE2025 V-C / IV-C (Calculation of dynamic score) + Algorithm 4
\subsection{動的評価スコア}
\label{subsec:holding_dynamic_score}

% --- 段落1 ---
動的評価スコア$S_{\mathrm{dyn}}$は，前節で生成された把持力候補群の中から，力量の大きさ・方向・分散の三基準に照らして最も安定と判定される力を把持点集合ごとに選出し，探索空間全体で正規化した評価値である（Algorithm 4）．

% --- 段落2 ---
動的評価は二段階で構成される．第一段階（局所評価）では，単一の把持点集合$P$に対して生成された全把持力候補を対象に，三基準の線形結合スコアを計算し，最高スコアに対応する力候補の素点を返す．

% --- 段落3 ---
力量の大きさ評価$e_{\mathrm{mag}}$は，全指先力のノルム和の負値として定義され，小さな力で把持可能な候補ほど高得点となる．これは，食材に加わる負荷が小さい方が望ましいという観点に基づく．

% --- 段落4 ---
方向評価$e_{\mathrm{dir}}$は，各指先力ベクトルと接触点における食材表面法線との内積和として定義される．指先力が表面法線方向に近いほど滑りにくいため，この値が大きい候補が選好される．

% --- 段落5 ---
分散評価$e_{\mathrm{var}}$は，各指先力の大きさの分散の負値として定義され，特定の指に力が集中することを避け，全指への均等な負荷分散を促す役割を担う．

% --- 段落6 ---
第二段階（大域評価）では，局所段階で選出された各把持点集合の素点を探索空間$\mathbb{P}$全体で$[0,1]$に正規化し，重み付き和として$S_{\mathrm{dyn}}$を算出する．本実験では，重みを$w_{\mathrm{dir}}=2,\ w_{\mathrm{mag}}=2,\ w_{\mathrm{var}}=1$と設定した．


% ------------------------------------------------------------
\section{動作制御：把持から切断までの実行}
\label{sec:holding_control}
% ------------------------------------------------------------

% ■ SICE2025 VI-A-1 (Hardwares) + VI-C (切断軌道と把持実行)
\subsection{把持点への手指配置と切断軌道の実行}
\label{subsec:holding_execution}

% --- 段落1 ---
実験プラットフォームの詳細は第3章にて述べた通りであり，左腕には食材形状への適応性を重視した3軸2指サーボハンドを，右腕には包丁を剛に固定するための高把持力平行グリッパを装着した（Fig. 3a, 3b）．

% --- 段落2 ---
探索により選定された把持点集合$P_{\mathrm{fin}}$の各接触位置へは，ヤコビアンに基づく指関節のみの逆運動学を用いて指先を誘導し把持を実行する．指先の位置制御にはPID制御を用いたが，サーボの駆動電流に設定された不感帯（最大電流の13\%）に起因して，指先位置に最大約1cmの誤差が生じうる点が実装上の制約として残されている．

% --- 段落3 ---
切断軌道としては，まな板面に対して25度の角度で下降する往復スライス軌道を用いた．この軌道は，第1相の定速前進下降と第2相の定速後進下降からなり，包丁が食材の初期接触高さの半分に達した時点で相を切り替える．切断面は，把持領域を十分に確保しかつ切断高さが3cm以上となるよう，各試行において人手により指定した．


% ------------------------------------------------------------
\section{実験と評価}
\label{sec:holding_experiments}
% ------------------------------------------------------------

% ■ SICE2025 VI-B (Experiment of generating holding force)
\subsection{把持力生成手法の効率性検証}
\label{subsec:holding_exp_force_generation}

% --- 段落1 ---
提案する把持力生成法の効率性を検証するため，6種類のジャガイモ点群データを用い，提案手法とランダム$\bm{k}$法のそれぞれについて，生成された把持力のうち摩擦円錐内に収まる割合（成功率）を比較した．

% --- 段落2 ---
実験では，各手法につき試行あたり$1.2\times 10^{5}$個の把持力を生成し，摩擦円錐の頂角は40度，軸方向は接触点における内向き法線ベクトルとした．また，幾何フィルタ後の探索空間サイズは上位1500件に固定した．

% --- 段落3 ---
実験結果をTable III（Table IIIの参照ラベル）に示す．提案手法の成功率は平均32.28\%であり，ランダム$\bm{k}$法の平均4.19\%に対して約7.7倍の生成効率を達成した．この結果から，提案手法が物理的に妥当な把持力候補をより効率的に生成できることが確かめられた．


% ■ SICE2025 VI-C (Experiment of real robot holding)
\subsection{把持安定性の比較評価}
\label{subsec:holding_exp_comparison}

% --- 段落1 ---
実ロボットによる切断実験では，6個の半切りジャガイモと4個の半切りタマネギを対象とし，提案手法（手法A）を含む4種類の把持点選択手法の安定性を比較した．

% --- 段落2 ---
手法BおよびCは，幾何フィルタ後の上位60\%候補の中から，それぞれ中位および最低位の評価スコアを持つ把持点集合を選択する条件である．手法Dは，幾何フィルタ適用前の全候補$\mathbb{P}_{\mathrm{all}}$からランダムに選択する条件である．

% --- 段落3 ---
各食材は4手法で1回ずつ切断され，切断前後に固定カメラで撮影した画像から食材の平行移動量（pixel）と回転量（deg）を計測することで把持安定性を評価した．なお，同一食材に対する繰り返し切断による切断力の低下を防ぐため，各切断後に計画切断面をその法線方向に2mmずつ移動させた．

% --- 段落4 ---
実験に用いた食材の物理パラメータは，Muら\cite{Mu2019RoboticCutting}の手法により測定した値を使用した．ジャガイモの破壊靭性は409.66\,J/m$^2$，圧力分布は3352.07\,N/m$^2$，タマネギの破壊靭性は646.46\,J/m$^2$，圧力分布は3971.2\,N/m$^2$である．評価関数の重みは，$w_{\mathrm{fin}}=1,\ w_{\mathrm{knf}}=4.4,\ w_{\mathrm{tbl}}=6$，$w_{\mathrm{pdir}}=5,\ w_{\mathrm{pdis}}=4$，$w_{\mathrm{dir}}=2,\ w_{\mathrm{mag}}=2,\ w_{\mathrm{var}}=1$とした．

% --- 段落5 ---
実験結果をTable IV（Table IVの参照ラベル）に示す．手法AおよびBは手法CおよびDと比較して明らかに小さな移動量を示し，評価スコアの高い把持点集合ほど安定した把持が実現できることが確認された．

% --- 段落6 ---
一方で，手法AとBの間には有意な差が見られなかった．これは，幾何フィルタ段階で明らかに不適切な候補が既に排除されており，上位候補はいずれも一定水準以上の安定性を有していたためと考察される．

% --- 段落7 ---
また，手法AおよびBにおいても一部の試行で大きな移動量（外れ値）が観測されたが，その主因はRGB-Dカメラ点群に残留するノイズによる形状モデルの誤差，およびサーボのPID制御における不感帯に起因する指先位置決め誤差（最大約1cm）であると特定された．

% --- 段落8 ---
さらに，本実験の評価法そのものにも限界がある．切断前後の画像比較のみでは切断途中の安定性を直接評価できず，往復スライス動作によって切断終了時に食材が元の位置近くに戻された場合，実際には途中で大きく移動していても変位が小さく評価される可能性がある．切断過程全体を通した把持安定性の評価は今後の課題である．


% ■ SICE2025 VI-D (Experiment of Complete Cutting)
\subsection{完全切断による実用性の実証}
\label{subsec:holding_exp_complete_cut}

% --- 段落1 ---
比較実験がサンプル形状保存のための部分切断にとどまるのに対し，本節では食材を完全に切断する実験を行い，提案手法の実用性を検証した．

% --- 段落2 ---
Fig. 5（Fig. 5の参照ラベル）に示す通り，提案手法により選定された把持点集合は，ジャガイモ・タマネギのいずれにおいても完全切断の全工程を通じて安定性を維持し，断面の滑らかな切断面が得られた．

% --- 段落3 ---
切断中に観測された変位は1cm未満であり，これも把持点での滑りではなく，サーボハンドの機械的剛性限界に起因するものであった．この結果は，提案手法が実用的な食材切断タスクに十分な把持性能を提供することを示している．


% ■ SICE2025 VI-E (Computation Time Analysis)
\subsection{計算コスト分析と実応用に向けた考察}
\label{subsec:holding_exp_computation}

% --- 段落1 ---
提案手法の実用展開を見据え，点群密度を7点から1220点まで変化させた際の計算時間と評価スコアの関係を分析した．全計算は，AMD Ryzen 9 3900X 12-Core Processor，64GB RAMを搭載した標準的なPC上で実行した．

% --- 段落2 ---
実験に採用した約170点の点群密度では平均計算時間は75.4秒であり，評価スコアは約22.85で飽和していることが確認された（Table II）．これ以上の点群密度増加はスコア改善にほとんど寄与せず，計算時間のみを増大させる．

% --- 段落3 ---
計算時間の大部分は動的評価スコアの計算（CalDynamicScore）が占めており，これは多数の把持力候補を生成・評価する処理に起因する．現状の実装はオフライン計画を前提としているが，家庭用ロボットが環境準備中やアプローチ動作中に把持計画を並行計算する用途には十分許容可能な遅延である．将来的には力評価処理のGPU並列化により，リアルタイム応用への展開も期待される．


% ------------------------------------------------------------
\section{本章のまとめ}
\label{sec:holding_summary}
% ------------------------------------------------------------

% --- 段落1 ---
本章では，包丁による食材切断時に生じる外力に対し，食材の意図しない移動を防ぐための把持点探索手法を提案した．

% --- 段落2 ---
提案手法の核心は，幾何的制約による候補空間の効率的削減，包丁が食材に与える力とトルクの力学的推定，把持力の冗長性を考慮した効率的な力候補生成，および大きさ・方向・分散の三基準に基づく動的評価スコアの統合にある．

% --- 段落3 ---
双腕ロボットシステムを用いた実機実験により，提案手法がランダム選択や低スコア候補と比較して有意に高い把持安定性を達成することが示された．さらに完全切断実験により，実用的な切断タスクへの適用可能性も確認された．

% --- 段落4 ---
一方で，現状の手法には平坦底面を前提とした運動制約，切断前後の画像比較に依存した評価法，および計算時間の制約という限界が残されている．底面が平坦でない食材や球状食材への拡張，切断過程全体を通した把持安定性の評価，および計算の高速化は今後の課題である．
